\documentclass[11pt]{article}
\newcommand{\videoid}{@VIDEOID@}
\input{../preambulo}
\curso{@CURSO@}
\unidad{@UNIDAD@}{@N@} % unidad @UNIDAD@ del curso; primer enunciado: @UNIDAD@.@N@

\title{@TITULO@}
\author{Arilín Haro\\[2pt]\small Facultad de Ciencias, UNAM}
\date{}

\begin{document}
\maketitle

\noindent\small Notas del video
\href{https://www.youtube.com/watch?v=@VIDEOID@}{\guillemotleft{}@TITULOVIDEO@\guillemotright{}}
(@CREDITO@curso de @NOMBRECURSO@). @RESUMEN@
\normalsize

% Contenido: seguir el video; no agregar resultados que no estén en él.
% Los comentarios propios se marcan: \paragraph{Comentario.} ...
% Enunciados numerados (un solo contador por unidad): \begin{defn}[Nombre], \begin{prop}, \begin{teo},
% \begin{lema}, \begin{cor}. Ejemplos sin número de unidad: \begin{ejem}.
\section*{Título de la primera sección}

% Solo si hay errores CONFIRMADOS por Arilín (ver notas/CONVENCIONES.md):
% \begin{erratas}
%   \errata{203}{3:23}{Se dice ..., pero debe decir ...}
% \end{erratas}

\end{document}
